\medskip

On note $f(t)$ la température (en $\degres C$) d'un café en fonction du temps $t$ (en minute) écoulé depuis sa sortie d'une machine à expresso. On admet, pour tout $t \in [0~;~ +\infty[$, l'expression suivante :

\[f(t) = 60\e^{- 0,08t} + 23.\]

Au bout de combien de temps la température du café sera-t-elle inférieure ou égale à $44~\degres C$ ?

Donner la réponse à la minute près.
